% -*- mode: latex; coding: utf-8 -*-
%-----------------------------------------------------------------------------
% pxLST -- User manual and API reference (truthful to repository state)
%
% Build from directory tex/:
%   set TEXINPUTS=.;./latex;./code;
%   pdflatex -halt-on-error -interaction=nonstopmode doc/pxLST.tex
%   pdflatex -halt-on-error -interaction=nonstopmode doc/pxLST.tex
%-----------------------------------------------------------------------------
\documentclass[11pt,a4paper]{article}
\usepackage[T1]{fontenc}
\usepackage[utf8]{inputenc}
\usepackage{lmodern}
\usepackage{hyperref}
\usepackage{pxLST}

\title{The \texttt{pxLST} package\\
  \large Professional code listings for P3CTeX}
\author{Scvria\\P3CTeX / UOC PEC Repository}
\date{2026/03/16 -- Version 0.1}

\begin{document}
\maketitle

\begin{abstract}
\texttt{pxLST} provides a small, stable \LaTeX2e\ API for typesetting
syntax-highlighted source code in styled \texttt{tcolorbox} frames using the
\texttt{listings} backend. The current implementation is fully contained in
\texttt{tex/latex/pxLST.sty}; \texttt{tex/code/pxLST.code.tex} is reserved for a
future split and is currently a stub.
\end{abstract}

\tableofcontents

\section{Loading}

\begin{verbatim}
\usepackage{pxLST}
\end{verbatim}

\section{Global configuration}

\subsection{\texttt{\textbackslash pxLSTsetup}}

\begin{verbatim}
\pxLSTsetup[<palette>]{<language>}[<lststyle>]
\end{verbatim}

\begin{itemize}
  \item \texttt{<palette>} is a palette name (default \texttt{dark}). The package
    ships \texttt{dark}, \texttt{pastel}, and \texttt{light}.
  \item \texttt{<language>} is a \texttt{listings} language name. The package
    defines: \texttt{pxJava}, \texttt{pxProlog}, \texttt{pxOWL}, \texttt{pxTAD},
    \texttt{pxYAML}, \texttt{algoritme}, \texttt{pxBash}, \texttt{pxPython},
    \texttt{pxC}, \texttt{pxJavaScript}, \texttt{pxJS}, \texttt{pxSQL}.
  \item \texttt{<lststyle>} is a \texttt{listings} style name (default
    \texttt{pxlst}). Additional styles provided: \texttt{block},
    \texttt{framed}, \texttt{plain}.
\end{itemize}

\paragraph{Example.}

\begin{verbatim}
\pxLSTsetup[dark]{pxJava}
\pxLSTsetup[pastel]{algoritme}[framed]
\pxLSTsetup[light]{pxYAML}
\end{verbatim}

\section{Environments}

All environments below are verbatim-safe and implemented via
\verb|\newtcblisting|. Palette/language are taken from the most recent
\verb|\pxLSTsetup| call.

\subsection{\texttt{Code} and \texttt{CODE}}

\begin{verbatim}
\begin{Code}[<listing options>]{<title>}
  ... verbatim code ...
\end{Code}
\end{verbatim}

\texttt{CODE} is an alias environment with the same argument contract.

\subsection{\texttt{miniCode}}

\begin{verbatim}
\begin{miniCode}[<width>]
  ... verbatim code ...
\end{miniCode}
\end{verbatim}

The default width is \verb|.8\linewidth|.

\subsection{\texttt{CodeBox} and \texttt{CodeBox*}}

\begin{verbatim}
\begin{CodeBox}[<listing options>]{<id>}{<title>}{<tag>}{<caption>}
  ... verbatim code ...
\end{CodeBox}
\end{verbatim}

\texttt{CodeBox*} has the same signature and is \emph{breakable} across pages.
Both environments insert \verb|\label{cdbx:<id>}| in the title bar, so you can
reference a box with \verb|\ref{cdbx:<id>}|.

\section{Commands}

\subsection{\texttt{\textbackslash pxLSTinput}}

\begin{verbatim}
\pxLSTinput[<language>]{<file>}
\end{verbatim}

Typesets an external file as a framed listing. The filename is resolved relative
to the working directory where \texttt{pdflatex} is executed.

\subsection{\texttt{\textbackslash pxCodeInline}}

\begin{verbatim}
\pxCodeInline[<language>]{<snippet>}
\end{verbatim}

Inline snippet rendered with \verb|\lstinline| in a small \texttt{tcolorbox}.
Content is \emph{not} verbatim (escape \verb|\%|, \verb|\#|, \verb|\&| as usual).

\section{Palette model}

Palettes are selected via the optional argument to \verb|\pxLSTsetup|:

\begin{verbatim}
\pxLSTsetup[dark]{pxJava}
\pxLSTsetup[pastel]{algoritme}
\pxLSTsetup[light]{pxYAML}
\end{verbatim}

The package provides three built-in palettes:

\begin{itemize}
  \item \texttt{dark}: high-contrast Darcula-inspired theme for on-screen use.
  \item \texttt{pastel}: softer variant of the dark theme.
  \item \texttt{light}: light background suitable for print.
\end{itemize}

Internally each palette is implemented by a macro that redefines a small family
of color-role commands to literal \texttt{xcolor} names (for example
\verb|\pxDEFlstfg| becomes \verb|\color{dark-fg}| in the dark palette). This keeps
the interaction with \texttt{listings} and \texttt{tcolorbox} predictable without
exposing additional configuration surface.

\end{document}
